\begin{methods}
\section{Methods}\label{sec:methods}

\subsection{Training data and labels}
100,000 random 9-mer peptides were generated by uniform sampling of the 20 standard amino acids. Each was labelled with MHCflurry 2.2.1 against HLA-A*02:01 and assigned to Strong Binder (SB, percentile rank $<$0.5\%), Weak Binder (WB, 0.5\%$\leq$rank$<$2.0\%), or Non-Binder (NB, rank $\geq$2.0\%) following the netMHCpan convention \citep{Reynisson20}. The resulting distribution (1.7\% SB, 5.1\% WB, 93.2\% NB) was balanced by downsampling to 1,730 per class (n=5,190), partitioned 90/10 into training (n=4,671) and held-out test (n=519) sets with stratified sampling (seed 42). Two comparator labelling strategies were implemented: a PSSM encoding crystallographic A*02:01 anchor preferences, and random synthetic labels as a baseline control. A held-out benchmark of 42 experimentally validated IEDB epitopes (HLA-A*02:01-restricted, with positive binding/T-cell assay evidence, re-verified against the IEDB on 2026-10-23) plus 20 homopolymer negatives ($\leq$2 unique amino acids) was curated from the IEDB (query-api.iedb.org); label-circularity overlap was determined by sequence-identity match against the IEDB entries in MHCflurry's training corpus \citep{Vita19}.

\subsection{Peptide encoding}
Each 9-mer was encoded with the BLOSUM62 matrix \citep{Henikoff92}, min--max normalized, producing a 9$\times$20 matrix flattened to 180 dimensions for fully connected networks. ESM-2 embeddings \citep{Lin23} were extracted at t6 (2,880-dim per-peptide; 320 per position $\times$ 9) and t12 (4,320-dim; 480 per position $\times$ 9) depths --- the two smallest ESM-2 checkpoints (8M and 35M parameters); a mean-pooled t6 variant (320-dim) served as an ablation control destroying positional information.

\subsection{Model architectures}
Six classifiers were implemented: (i) FFN \citep{Nielsen03}: 180$\rightarrow$90$\rightarrow$3, 49K params; (ii) Deep FFN: 180$\rightarrow$360$\rightarrow$180$\rightarrow$90$\rightarrow$45$\rightarrow$3 with batch normalisation and dropout, 152K params; (iii) CNN: Conv2D(32,3$\times$3) + FFN body, 1.05M params; (iv) LSTM(64)$\rightarrow$dense(32)$\rightarrow$3, 24K params; (v) ResNet-style stem + 3 residual blocks, 325K params; (vi) Random Forest (100 trees). For ESM-2 embeddings an FFN classifier with a 256$\rightarrow$128$\rightarrow$64 stack was trained. All networks used cross-entropy loss with Adam or RMSprop (lr 0.001), 150 epochs, batch size 50, 20\% internal validation split, early stopping (patience 10), and learning-rate reduction on plateau. Five-fold stratified cross-validation was performed for the FFN. The held-out test set (519 peptides) was disjoint from training and used for all reported metrics. Architecture comparisons were repeated across five fixed training seeds; per-seed variation reflects data-sampling randomness, and seeds are not independent replicates.

\subsection{Protein scanning, mutation analysis, and structure}
Ten therapeutically relevant proteins (MART-1, gp100, tyrosinase, NY-ESO-1, WT1, p53, KRAS, influenza M1, CMV pp65, SARS-CoV-2 spike RBD) were scanned with a sliding 9-mer window (stride 1; 3,536 windows). Sixteen recurrent hotspot mutations (7 p53, 9 KRAS) were analyzed by comparing wild-type and mutant binding scores; mutations were classified as Created/Enhanced/Destroyed/Unchanged. The KRAS G12V neoepitope YKLVVVGAV carries a non-canonical lysine at P2.

\subsection{Feature extraction}
Learned features were extracted from the Deep FFN penultimate dense layer via truncated models. Hierarchical clustering (Ward's method, Euclidean distance) was applied to representations of top epitope candidates; feature--binding-score correlation was assessed by Pearson's $r$ with a 1,000-iteration permutation test to control for chance correlations. t-SNE visualisation of the 45-dimensional penultimate layer was used to inspect class separation, and integrated gradients identified position-specific anchor contributions.

\subsection{Multi-allele extension}
\subsubsection{Multi-allele data curation}
Experimentally confirmed class I binders were curated from the IEDB \citep{Vita19} for five common alleles (HLA-A*02:01, A*01:01, A*03:01, B*07:02, B*44:03), yielding n = 14,488 positives (canonical 11,225; Type A 151; Type B 3,112). A 15-allele panel (pan16: 11 additional alleles --- A*11:01, A*24:02, B*08:01, B*15:01, B*27:05, B*35:01, B*40:01, B*49:01, B*51:01, B*57:01) was curated for generalisation testing. Protein-background negatives were sampled from natural proteomes to avoid the uniform-random decoys used in earlier evaluations. IEDB data were queried live from query-api.iedb.org (accessed 2026-08-02); the 15-allele and temporal-holdout estimates are therefore snapshot-dependent and drift slightly as the IEDB grows.

\subsubsection{Type A/B classification}
Type A (tolerated weak anchors) and Type B (energy-compensated non-canonical) were defined by whether the peptide violates the allele's canonical anchor motif at p2/p9. Type A carries broad-only tolerated anchors (residues that appear in binding motifs but are non-canonical); Type B carries non-canonical residues at anchor positions. The two classes have different biological origins and, as we show, different prediction difficulty.

\subsubsection{Pan model and per-allele comparators}
The pan-specific DFFN uses BLOSUM62 encoding \citep{Henikoff92} through a 360--180--90--45 hidden stack with batch normalisation and dropout. Rebalancing upsampled the minority B*07:02 allele 4$\times$ during training to correct for allele-frequency imbalance. Five seeds (0, 1, 42, 123, 2024) were used; per-seed variation arises from data-sampling RNG. Per-allele models (one per allele) and a pan16 model (15 alleles, per-allele minimum 2,000 upsampling) were trained with the same recipe.

\subsubsection{Comparator predictors}
MHCflurry 2.2.1 \citep{ODonnell20} served as an independent sequence predictor (affinity threshold 500 nM). MUNIS \citep{Wohlwend25} --- a second sequence-based predictor using ESM-2 embeddings \citep{Rives21} --- served as an independent sequence comparator, run on the n = 3,112 Type B population plus canonical controls.

\subsubsection{Temporal holdout}
To address label circularity \citep{Preibisch26}, IEDB epitopes published after 2022 (post-2022, unseen by MHCflurry's $\sim$pre-2021 training data) were compared against pre-2021 (in-distribution) epitopes at sensitivity threshold 0.95.

\subsubsection{Evaluation}
Sensitivity was computed with protein-background negatives and fixed thresholds (0.95 for the pan model; 500 nM for MHCflurry; 5th percentile of canonical scores for MUNIS and BigMHC). Wilson 95\% confidence intervals are reported for pooled estimates.

\end{methods}
